% \section{Identifiability Theory for Source Apportionment}
% \label{sec:identifiability_theory}

% Section~\ref{sec:wind_conditioned_source_apportionment} constructs a
% lagged wind-conditioned inverse problem for estimating source activity.
% However, a good sensor fit alone does not imply that the inferred source
% contributions are uniquely determined by the observations. In this
% section, we characterize when inventory-based source apportionment is
% identifiable from sparse sensor data.

% The central object is the projected lagged source-response matrix
% \begin{equation}
% \widetilde H
% :=
% P_Q^\perp H_{1:T}^{\mathrm{lag}}
% \in
% \mathbb R^{mT\times K},
% \qquad
% \widetilde{\mathbf Y}
% :=
% P_Q^\perp \mathbf Y
% \in
% \mathbb R^{mT}.
% \label{eq:theory_projected_matrix}
% \end{equation}
% The projected observation model is
% \begin{equation}
% \widetilde{\mathbf Y}
% =
% \widetilde H\boldsymbol\theta
% +
% \widetilde{\mathbf E},
% \qquad
% \widetilde{\mathbf E}:=P_Q^\perp\mathbf E.
% \label{eq:projected_identifiability_model}
% \end{equation}
% Thus, identifiability is governed not by the raw response matrix
% \(H_{1:T}^{\mathrm{lag}}\), but by the response that remains after
% background and temporal correction.

% \subsection{Exact Identifiability}

% We first consider the noiseless case. The source contribution vector
% \(\boldsymbol\theta\) is exactly identifiable if no other source
% contribution vector produces the same projected sensor trajectory.

% \begin{definition}[Exact identifiability]
% The source contribution vector \(\boldsymbol\theta\) is identifiable from
% \(\widetilde{\mathbf Y}\) if
% \begin{equation}
% \widetilde H\boldsymbol\theta_1
% =
% \widetilde H\boldsymbol\theta_2
% \quad
% \Longrightarrow
% \quad
% \boldsymbol\theta_1=\boldsymbol\theta_2 .
% \end{equation}
% \end{definition}

% \begin{proposition}[Identifiability of inventory-based apportionment]
% \label{prop:exact_identifiability}
% In the noiseless projected model
% \begin{equation}
% \widetilde{\mathbf Y}
% =
% \widetilde H\boldsymbol\theta,
% \end{equation}
% the source contribution vector is identifiable for all
% \(\boldsymbol\theta\in\mathbb R^K\) if and only if
% \begin{equation}
% \operatorname{rank}(\widetilde H)=K.
% \label{eq:full_rank_condition}
% \end{equation}
% \end{proposition}

% \begin{proof}
% If \(\operatorname{rank}(\widetilde H)=K\), then
% \(\ker(\widetilde H)=\{\mathbf 0\}\). Therefore
% \(\widetilde H\boldsymbol\theta_1=
% \widetilde H\boldsymbol\theta_2\) implies
% \(\widetilde H(\boldsymbol\theta_1-\boldsymbol\theta_2)=0\), and hence
% \(\boldsymbol\theta_1=\boldsymbol\theta_2\).

% Conversely, if \(\operatorname{rank}(\widetilde H)<K\), then there exists
% a nonzero vector \(\Delta\boldsymbol\theta\in\ker(\widetilde H)\). Hence
% \[
% \widetilde H\boldsymbol\theta
% =
% \widetilde H(\boldsymbol\theta+\Delta\boldsymbol\theta),
% \]
% so the source contribution vector is not uniquely determined.
% \end{proof}

% The full-rank condition is the global linear identifiability condition.
% Since source activities are constrained to be nonnegative in our
% apportionment model, nonnegativity can sometimes remove some ambiguous
% directions. However, full column rank remains the cleanest sufficient
% condition for stable source-level interpretation. If
% \(\operatorname{rank}(\widetilde H)<K\), then any source activity vector
% in the interior of the nonnegative cone can be perturbed slightly along a
% null direction while remaining nonnegative, producing the same projected
% observations.

% \subsection{Noise-Robust Identifiability}

% In practice, observations contain sensor noise, wind interpolation error,
% source-inventory uncertainty, and transport-model approximation error.
% Even when \(\widetilde H\) has full rank, recovery may be unstable if its
% columns are nearly linearly dependent.

% Let
% \[
% \sigma_1(\widetilde H)
% \ge
% \sigma_2(\widetilde H)
% \ge
% \cdots
% \ge
% \sigma_K(\widetilde H)
% \]
% denote the singular values of \(\widetilde H\). The smallest singular
% value controls the worst-case amplification of observation error.

% \begin{proposition}[Noise-robust apportionment]
% \label{prop:noise_robust_apportionment}
% Assume
% \[
% \widetilde{\mathbf Y}
% =
% \widetilde H\boldsymbol\theta
% +
% \widetilde{\mathbf E}
% \]
% and \(\operatorname{rank}(\widetilde H)=K\). Let
% \(\widehat{\boldsymbol\theta}_{\mathrm{LS}}\) be the unconstrained
% least-squares estimate. Then
% \begin{equation}
% \left\|
% \widehat{\boldsymbol\theta}_{\mathrm{LS}}
% -
% \boldsymbol\theta
% \right\|_2
% \le
% \frac{
% \|\widetilde{\mathbf E}\|_2
% }{
% \sigma_{\min}(\widetilde H)
% }.
% \label{eq:ls_noise_bound}
% \end{equation}
% If \(\widehat{\boldsymbol\theta}\) is the nonnegative least-squares
% estimate and the true \(\boldsymbol\theta\ge 0\), then
% \begin{equation}
% \left\|
% \widehat{\boldsymbol\theta}
% -
% \boldsymbol\theta
% \right\|_2
% \le
% \frac{
% 2\|\widetilde{\mathbf E}\|_2
% }{
% \sigma_{\min}(\widetilde H)
% }.
% \label{eq:nnls_noise_bound}
% \end{equation}
% \end{proposition}

% \begin{proof}
% For the unconstrained least-squares estimator,
% \[
% \widehat{\boldsymbol\theta}_{\mathrm{LS}}
% -
% \boldsymbol\theta
% =
% (\widetilde H^\top \widetilde H)^{-1}
% \widetilde H^\top \widetilde{\mathbf E}.
% \]
% Taking norms gives
% \[
% \left\|
% \widehat{\boldsymbol\theta}_{\mathrm{LS}}
% -
% \boldsymbol\theta
% \right\|_2
% \le
% \|\widetilde H^\dagger\|_2
% \|\widetilde{\mathbf E}\|_2
% =
% \frac{\|\widetilde{\mathbf E}\|_2}
% {\sigma_{\min}(\widetilde H)}.
% \]

% For nonnegative least squares, the true
% \(\boldsymbol\theta\ge 0\) is feasible. By optimality,
% \[
% \|\widetilde{\mathbf Y}-\widetilde H\widehat{\boldsymbol\theta}\|_2
% \le
% \|\widetilde{\mathbf Y}-\widetilde H\boldsymbol\theta\|_2
% =
% \|\widetilde{\mathbf E}\|_2.
% \]
% Therefore
% \[
% \|\widetilde H(\widehat{\boldsymbol\theta}-\boldsymbol\theta)\|_2
% \le
% 2\|\widetilde{\mathbf E}\|_2.
% \]
% Since \(\widetilde H\) has full column rank,
% \[
% \|\widetilde H(\widehat{\boldsymbol\theta}-\boldsymbol\theta)\|_2
% \ge
% \sigma_{\min}(\widetilde H)
% \|\widehat{\boldsymbol\theta}-\boldsymbol\theta\|_2,
% \]
% which proves the claim.
% \end{proof}

% This result shows that source apportionment may be formally identifiable
% but practically unreliable when \(\sigma_{\min}(\widetilde H)\) is small.

% \subsection{Identifiability Scores}

% \paragraph{Identifiability Score} We define the primary identifiability score as
% \begin{equation}
% \mathcal I
% :=
% \sigma_{\min}(\widetilde H).
% \label{eq:identifiability_score}
% \end{equation}
% Larger \(\mathcal I\) implies stronger separation between source
% fingerprints and lower worst-case sensitivity to noise.

% \paragraph{Condition Number} We also define the condition number
% \begin{equation}
% \kappa(\widetilde H)
% =
% \frac{
% \sigma_{\max}(\widetilde H)
% }{
% \sigma_{\min}(\widetilde H)
% }.
% \label{eq:condition_number}
% \end{equation}
% Large \(\kappa(\widetilde H)\) indicates that some linear combinations of
% source groups are much harder to estimate than others.

% \paragraph{Effective Rank} For a noise-dependent threshold \(\tau_\sigma>0\), we define the
% effective rank
% \begin{equation}
% r_{\mathrm{eff}}(\tau_\sigma)
% =
% \#\left\{
% i:
% \sigma_i(\widetilde H)>\tau_\sigma
% \right\}.
% \label{eq:effective_rank}
% \end{equation}
% If
% \begin{equation}
% r_{\mathrm{eff}}(\tau_\sigma)<K,
% \end{equation}
% then the sensor network does not support reliable estimation of all
% \(K\) source groups at the given noise level.

% \paragraph{Source Visibility}

% Rank and conditioning describe global identifiability. We also need
% source-specific diagnostics. Let
% \[
% \widetilde{\mathbf h}_k
% :=
% P_Q^\perp
% \mathbf h_k^{\mathrm{lag}}
% \]
% denote the projected lagged fingerprint of source group \(k\), where
% \(\mathbf h_k^{\mathrm{lag}}\) is the \(k\)-th column of
% \(H_{1:T}^{\mathrm{lag}}\). We define the visibility of source group
% \(k\) as
% \begin{equation}
% v_k
% :=
% \|\widetilde{\mathbf h}_k\|_2.
% \label{eq:source_visibility}
% \end{equation}
% A small value of \(v_k\) means that source group \(k\) has little
% measurable effect on the sensor network after lagged transport and
% background correction. Such a source may be physically present in the
% inventory but weakly observed by the sensing geometry and wind regime.

% % We can also define normalized visibility
% % \begin{equation}
% % \bar v_k
% % =
% % \frac{
% % \|\widetilde{\mathbf h}_k\|_2
% % }{
% % \|\mathbf s_k\|_2
% % +
% % \varepsilon
% % },
% % \label{eq:normalized_visibility}
% % \end{equation}
% % which measures observable sensor response per unit inventory magnitude.

% \paragraph{Pairwise Source Ambiguity}

% Two source groups may both be visible but still difficult to separate if
% their projected fingerprints are nearly proportional. We quantify this
% using projected source-fingerprint coherence:
% \begin{equation}
% \rho_{ij}
% =
% \frac{
% \left|
% \widetilde{\mathbf h}_i^\top
% \widetilde{\mathbf h}_j
% \right|
% }{
% \|\widetilde{\mathbf h}_i\|_2
% \|\widetilde{\mathbf h}_j\|_2
% }.
% \label{eq:projected_fingerprint_coherence}
% \end{equation}
% Values \(\rho_{ij}\approx 1\) indicate that source groups \(i\) and \(j\)
% produce nearly collinear sensor-time signatures after background
% correction.

% This captures a common ambiguity in sparse air-quality apportionment: a
% strong distant source and a weak nearby source can generate similar
% sensor readings if their lagged wind-conditioned fingerprints are nearly
% proportional. In this case, the sensor data may constrain the combined
% effect of the sources but not their separate contributions.

% % To express this scale-invariant ambiguity directly, define the ray
% % distance
% % \begin{equation}
% % d^{\mathrm{ray}}_{ij}
% % =
% % \min_{c\in\mathbb R}
% % \frac{
% % \left\|
% % \widetilde{\mathbf h}_i
% % -
% % c\widetilde{\mathbf h}_j
% % \right\|_2
% % }{
% % \|\widetilde{\mathbf h}_i\|_2
% % }.
% % \label{eq:ray_distance}
% % \end{equation}
% % For nonzero projected fingerprints, this satisfies
% % \begin{equation}
% % d^{\mathrm{ray}}_{ij}
% % =
% % \sqrt{1-\rho_{ij}^2}.
% % \label{eq:ray_distance_coherence_relation}
% % \end{equation}
% % Thus, high coherence is equivalent to small ray distance.

% \subsection{Background Correction Can Reduce Identifiability}

% The background-correction projection changes the source-response matrix.
% A source group that appears visible in the raw response matrix may become
% weakly visible after projection if its fingerprint lies close to the
% background-correction subspace.

% For source group \(k\), the removed fraction is
% \begin{equation}
% a_k
% =
% \frac{
% \|P_Q\mathbf h_k^{\mathrm{lag}}\|_2
% }{
% \|\mathbf h_k^{\mathrm{lag}}\|_2
% },
% \qquad
% P_Q:=QQ^\dagger.
% \label{eq:background_absorption_fraction}
% \end{equation}
% If \(a_k\approx 1\), then most of the source fingerprint can be explained
% by the background-correction basis. The projected visibility
% \(\|\widetilde{\mathbf h}_k\|_2\) will be small even if the raw
% fingerprint is large.

% This is why identifiability must be assessed using
% \(\widetilde H=P_Q^\perp H_{1:T}^{\mathrm{lag}}\), not the unprojected
% matrix \(H_{1:T}^{\mathrm{lag}}\).

% \subsection{Operator Error and Wind Uncertainty}

% The response matrix itself may be imperfect because wind fields,
% dispersion parameters, source inventories, and plume approximations are
% estimated. Let \(H^\star\) be the true lagged response matrix and
% \(\widehat H\) the constructed response matrix. After background
% projection, define
% \[
% \widetilde H^\star=P_Q^\perp H^\star,
% \qquad
% \widehat{\widetilde H}=P_Q^\perp \widehat H.
% \]
% Suppose
% \begin{equation}
% \left\|
% \widehat{\widetilde H}
% -
% \widetilde H^\star
% \right\|_2
% \le
% \delta_H.
% \label{eq:operator_error_bound}
% \end{equation}
% Then the projected observations satisfy
% \begin{equation}
% \widetilde{\mathbf Y}
% =
% \widehat{\widetilde H}\boldsymbol\theta
% +
% \left[
% \widetilde{\mathbf E}
% +
% (\widetilde H^\star-\widehat{\widetilde H})
% \boldsymbol\theta
% \right].
% \label{eq:operator_error_as_noise}
% \end{equation}
% Thus, response-matrix error acts like an additional observation error.

% \begin{proposition}[Error bound under response-matrix perturbation]
% \label{prop:operator_error_bound}
% Assume \(\widehat{\widetilde H}\) has full column rank and
% \[
% \left\|
% \widehat{\widetilde H}
% -
% \widetilde H^\star
% \right\|_2
% \le
% \delta_H.
% \]
% Then the unconstrained least-squares estimate using
% \(\widehat{\widetilde H}\) satisfies
% \begin{equation}
% \left\|
% \widehat{\boldsymbol\theta}
% -
% \boldsymbol\theta
% \right\|_2
% \le
% \frac{
% \|\widetilde{\mathbf E}\|_2
% +
% \delta_H\|\boldsymbol\theta\|_2
% }{
% \sigma_{\min}(\widehat{\widetilde H})
% }.
% \label{eq:operator_error_theta_bound}
% \end{equation}
% \end{proposition}

% This bound shows that wind interpolation uncertainty and plume-model
% error reduce attribution reliability in the same way as sensor noise:
% their effect is amplified when the projected response matrix is
% ill-conditioned.

% \subsection{Identifiable Source Resolution and Merging}

% When source fingerprints are weak or highly coherent, reporting separate
% source contributions can be misleading. Instead, the source partition
% should be coarsened until the projected response matrix is sufficiently
% well-conditioned.

% Let \(\mathcal G=\{G_1,\ldots,G_M\}\) be a partition of the \(K\) source
% groups into \(M\) merged groups. The merged inventory matrix is
% \begin{equation}
% S^{\mathcal G}
% =
% \begin{bmatrix}
% \sum_{k\in G_1}\mathbf s_k &
% \sum_{k\in G_2}\mathbf s_k &
% \cdots &
% \sum_{k\in G_M}\mathbf s_k
% \end{bmatrix}.
% \label{eq:merged_inventory_matrix}
% \end{equation}
% Using \(S^{\mathcal G}\), we construct a merged projected response matrix
% \(\widetilde H^{\mathcal G}\).

% We say that a source partition is identifiable at threshold
% \(\tau_\sigma\) if
% \begin{equation}
% r_{\mathrm{eff}}
% \left(
% \widetilde H^{\mathcal G};
% \tau_\sigma
% \right)
% =
% M
% \label{eq:identifiable_partition_condition}
% \end{equation}
% and
% \begin{equation}
% \max_{a\neq b}
% \rho_{ab}^{\mathcal G}
% \le
% \tau_\rho,
% \label{eq:merged_coherence_condition}
% \end{equation}
% where \(\rho_{ab}^{\mathcal G}\) is the coherence between merged group
% fingerprints and \(\tau_\rho\in(0,1)\) is a user-specified ambiguity
% threshold.

% This gives a practical reporting rule:

% \begin{quote}
% Report source contributions at the finest source grouping whose projected
% lagged response matrix is sufficiently well-conditioned and whose
% pairwise source-fingerprint coherence remains below the ambiguity
% threshold.
% \end{quote}

% If two source groups have \(\rho_{ij}>\tau_\rho\), they should be treated
% as ambiguous under the current sensor layout, wind regime, and background
% correction. The combined contribution may be identifiable even when the
% separate contributions are not.

% \subsection{Summary of Identifiability Diagnostics}

% The theory suggests that an apportionment method should report the
% following diagnostics in addition to source contribution estimates:
% \begin{enumerate}[leftmargin=*]
%     \item rank of \(\widetilde H\);
%     \item smallest singular value
%     \(\sigma_{\min}(\widetilde H)\);
%     \item condition number \(\kappa(\widetilde H)\);
%     \item effective rank \(r_{\mathrm{eff}}(\tau_\sigma)\);
%     \item source visibility \(v_k\);
%     \item background absorption fraction \(a_k\);
%     \item pairwise coherence \(\rho_{ij}\);
%     % \item ray distance \(d^{\mathrm{ray}}_{ij}\);
%     \item source groups recommended for merging.
% \end{enumerate}

% Together, these quantities define the identifiable resolution of the
% source apportionment problem. Sparse sensors may support reliable
% estimation of broad source categories while failing to support
% fine-grained separation between highly coherent source groups. The role
% of the method is therefore not only to estimate
% \(\boldsymbol\theta\), but also to report which components of
% \(\boldsymbol\theta\) are defensible from the available observations.



\section{Identifiability Theory for Source Apportionment}
\label{sec:identifiability_theory}

We now characterize when the source activity vector is determined by the
projected observations.  Throughout this section,
\begin{equation}
\widetilde Y=\widetilde H\boldsymbol\theta+\widetilde E,
\qquad
\widetilde H=P_Q^\perp H_{1:T}^{\mathrm{lag}}\in\mathbb R^{mT\times K}.
\label{eq:theory_projected_model}
\end{equation}
The analysis is conditional on the constructed wind field, transport operator,
source inventories, lag window, and background basis.

\subsection{Exact Identifiability}
\label{subsec:exact_identifiability}

\begin{definition}[Exact identifiability]
The source activity vector is identifiable from \(\widetilde Y\) if
\begin{equation}
\widetilde H\boldsymbol\theta_1=\widetilde H\boldsymbol\theta_2
\quad\Longrightarrow\quad
\boldsymbol\theta_1=\boldsymbol\theta_2.
\label{eq:exact_identifiability_definition}
\end{equation}
\end{definition}

\begin{proposition}[Identifiability of inventory-based apportionment]
\label{prop:rank_identifiability}
In the noiseless projected model \(\widetilde Y=\widetilde H\boldsymbol\theta\),
\(\boldsymbol\theta\in\mathbb R^K\) is identifiable for all source activity
vectors if and only if
\begin{equation}
\operatorname{rank}(\widetilde H)=K.
\label{eq:rank_condition}
\end{equation}
\end{proposition}

\begin{proof}
If \(\operatorname{rank}(\widetilde H)=K\), then
\(\ker(\widetilde H)=\{0\}\), so
\(\widetilde H(\boldsymbol\theta_1-\boldsymbol\theta_2)=0\) implies
\(\boldsymbol\theta_1=\boldsymbol\theta_2\).  Conversely, if
\(\operatorname{rank}(\widetilde H)<K\), there exists a nonzero
\(\Delta\boldsymbol\theta\in\ker(\widetilde H)\).  Then
\(\widetilde H\boldsymbol\theta=\widetilde H(\boldsymbol\theta+
\Delta\boldsymbol\theta)\), so the source vector is not uniquely determined.
\end{proof}

Nonnegativity can remove some null directions at the boundary of the feasible
cone.  However, if \(\operatorname{rank}(\widetilde H)<K\), any activity vector
in the interior of \(\mathbb R_+^K\) can be perturbed slightly along a null
direction while remaining feasible.  Full column rank is therefore the clean
condition for source-level interpretation.

\subsection{Noise-Robust Identifiability}
\label{subsec:noise_robust_identifiability}

Let \(\sigma_1(\widetilde H)\ge\cdots\ge\sigma_K(\widetilde H)\) be the singular
values of \(\widetilde H\).  Even with full rank, small singular values amplify
noise.

\begin{proposition}[Noise-robust apportionment]
\label{prop:noise_robust}
Assume \(\widetilde Y=\widetilde H\boldsymbol\theta+\widetilde E\) and
\(\operatorname{rank}(\widetilde H)=K\).  Let
\(\widehat{\boldsymbol\theta}_{\mathrm{LS}}\) be the unconstrained least-squares
estimate.  Then
\begin{equation}
\|\widehat{\boldsymbol\theta}_{\mathrm{LS}}-\boldsymbol\theta\|_2
\le
\frac{\|\widetilde E\|_2}{\sigma_{\min}(\widetilde H)}.
\label{eq:ls_noise_bound}
\end{equation}
If \(\widehat{\boldsymbol\theta}\) is the nonnegative least-squares estimate and
\(\boldsymbol\theta\ge0\), then
\begin{equation}
\|\widehat{\boldsymbol\theta}-\boldsymbol\theta\|_2
\le
\frac{2\|\widetilde E\|_2}{\sigma_{\min}(\widetilde H)}.
\label{eq:nnls_noise_bound}
\end{equation}
\end{proposition}

\begin{proof}
For least squares,
\(\widehat{\boldsymbol\theta}_{\mathrm{LS}}-\boldsymbol\theta=
\widetilde H^\dagger\widetilde E\), giving
\(\|\widehat{\boldsymbol\theta}_{\mathrm{LS}}-\boldsymbol\theta\|_2\le
\|\widetilde H^\dagger\|_2\|\widetilde E\|_2=
\|\widetilde E\|_2/\sigma_{\min}(\widetilde H)\).  For nonnegative least
squares, the true \(\boldsymbol\theta\) is feasible, so optimality gives
\(\|\widetilde Y-\widetilde H\widehat{\boldsymbol\theta}\|_2\le
\|\widetilde E\|_2\).  Hence
\(\|\widetilde H(\widehat{\boldsymbol\theta}-\boldsymbol\theta)\|_2\le
2\|\widetilde E\|_2\).  The lower singular-value bound completes the proof.
\end{proof}

Thus \(\sigma_{\min}(\widetilde H)\) measures worst-case robustness: source
apportionment can be exactly identifiable but practically unstable when this
quantity is small.

\subsection{Identifiability Diagnostics}
\label{subsec:identifiability_diagnostics}

The primary identifiability score is
\begin{equation}
\mathcal I = \sigma_{\min}(\widetilde H).
\label{eq:identifiability_score}
\end{equation}
We also use the condition number
\begin{equation}
\kappa(\widetilde H)=
\frac{\sigma_{\max}(\widetilde H)}{\sigma_{\min}(\widetilde H)},
\label{eq:condition_number}
\end{equation}
and, for a noise-dependent threshold \(\tau_\sigma\), the effective rank
\begin{equation}
r_{\mathrm{eff}}(\tau_\sigma)
=\#\{i:\sigma_i(\widetilde H)>\tau_\sigma\}.
\label{eq:effective_rank}
\end{equation}
If \(r_{\mathrm{eff}}(\tau_\sigma)<K\), the sensor network does not support
reliable estimation of all \(K\) source groups at that noise level.

For source-specific diagnostics, let \(\widetilde{\mathbf h}_k\) be the
projected fingerprint of source group \(k\).  Its visibility is
\begin{equation}
v_k=\|\widetilde{\mathbf h}_k\|_2.
\label{eq:visibility}
\end{equation}
A small \(v_k\) means that the source has little measurable effect after
transport and background correction.  Pairwise ambiguity is measured by
projected fingerprint coherence,
\begin{equation}
\rho_{ij}
=
\frac{|\widetilde{\mathbf h}_i^\top\widetilde{\mathbf h}_j|}
{\|\widetilde{\mathbf h}_i\|_2\|\widetilde{\mathbf h}_j\|_2}.
\label{eq:pairwise_coherence}
\end{equation}
Values near one indicate that two source groups produce nearly proportional
sensor-time signatures.  In that case, the data may constrain their combined
contribution better than their separate contributions.

\subsection{Background Absorption}
\label{subsec:background_absorption}

Background correction can improve robustness by removing broad temporal
variation, but it can also remove source-driven signal.  Let \(P_Q=QQ^\dagger\).
For source group \(k\), define the background absorption fraction
\begin{equation}
a_k=
\frac{\|P_Q\mathbf h_k^{\mathrm{lag}}\|_2}
{\|\mathbf h_k^{\mathrm{lag}}\|_2}.
\label{eq:background_absorption}
\end{equation}
If \(a_k\approx 1\), most of the raw source fingerprint lies in the background
space, and the projected visibility \(v_k\) will be small even if the raw
fingerprint is large.  This is why identifiability must be assessed after
projection.

\subsection{Response-Matrix Perturbations}
\label{subsec:operator_error}

The constructed response matrix can be wrong because of wind interpolation,
transport approximation, dispersion parameters, or inventory error.  Let
\(\widetilde H^\star\) be the true projected response and
\(\widehat{\widetilde H}\) the constructed one, with
\begin{equation}
\|\widehat{\widetilde H}-\widetilde H^\star\|_2\le\delta_H.
\label{eq:operator_perturbation}
\end{equation}
Then
\begin{equation}
\widetilde Y
=\widehat{\widetilde H}\boldsymbol\theta
+\left[\widetilde E+(\widetilde H^\star-\widehat{\widetilde H})\boldsymbol\theta\right].
\label{eq:operator_error_as_noise}
\end{equation}
Response error therefore acts like additional observation error.

\begin{proposition}[Perturbation bound]
\label{prop:operator_perturbation}
Assume \(\widehat{\widetilde H}\) has full column rank and
\(\|\widehat{\widetilde H}-\widetilde H^\star\|_2\le\delta_H\).  The
unconstrained least-squares estimate using \(\widehat{\widetilde H}\) satisfies
\begin{equation}
\|\widehat{\boldsymbol\theta}-\boldsymbol\theta\|_2
\le
\frac{\|\widetilde E\|_2+\delta_H\|\boldsymbol\theta\|_2}
{\sigma_{\min}(\widehat{\widetilde H})}.
\label{eq:operator_error_bound}
\end{equation}
\end{proposition}

The effect of wind and transport uncertainty is therefore most damaging when
the projected response matrix is already ill-conditioned.

\subsection{Identifiable Source Resolution}
\label{subsec:source_merging}

When source fingerprints are weak or highly coherent, separate estimates can be
misleading.  Let \(\mathcal G=\{G_1,\ldots,G_M\}\) be a partition of the original
\(K\) source groups into merged groups, with merged inventory
\begin{equation}
S^{\mathcal G}
=
\left[
\sum_{k\in G_1}\mathbf s_k\;
\sum_{k\in G_2}\mathbf s_k\;
\cdots\;
\sum_{k\in G_M}\mathbf s_k
\right].
\label{eq:merged_inventory}
\end{equation}
Let \(\widetilde H^{\mathcal G}\) be the corresponding projected response.  We
say that \(\mathcal G\) is identifiable at thresholds \(\tau_\sigma\) and
\(\tau_\rho\) if
\begin{equation}
r_{\mathrm{eff}}(\widetilde H^{\mathcal G};\tau_\sigma)=M,
\qquad
\max_{a\ne b}\rho^{\mathcal G}_{ab}\le \tau_\rho.
\label{eq:merged_identifiability}
\end{equation}
The reporting rule is to use the finest partition satisfying these conditions,
while flagging low-visibility groups.

\begin{table}[t]
\centering
\small
\caption{Identifiability diagnostics reported with source estimates.}
\label{tab:diagnostics}
\begin{tabular}{lll}
\toprule
Diagnostic & Interpretation & Action \\
\midrule
\(\operatorname{rank}(\widetilde H)\) & exact separability & require full rank \\
\(\sigma_{\min}(\widetilde H)\) & noise robustness & flag instability \\
\(\kappa(\widetilde H)\) & error amplification & compare settings \\
\(r_{\mathrm{eff}}\) & noise-level resolution & reduce grouping \\
\(v_k\) & source visibility & flag weak sources \\
\(a_k\) & background absorption & warn after projection \\
\(\rho_{ij}\) & pairwise ambiguity & merge coherent groups \\
\bottomrule
\end{tabular}
\end{table}
